% topics/polytope-geometry/exact-polar-geometry.tex — ordinary exact normalized H-geometry
% Imported once by formal/main.tex. Labels moved from
% symplectic-polytope-geometry.tex, random-polytope-boundedness.tex and
% rational-integer-scaling.tex retain their identifiers.
% Provenance: the boundedness proposition and original proof were
% Jörn-reviewed (57e73d9). Its interval-endpoint argument below is corrected
% in this consolidation and awaits Jörn's review. The other agent-written
% geometry and scaling notes remain unverified unless independently marked.

\subsection{Normalized polar geometry}

\begin{unverified}
\begin{definition}[Polytope dual / polar body]\label{def:polytope-dual}
% [TODO: JÖRN - verify statement]
% Source: polytope.rs doc comment
Let $K \subset \mathbb{R}^4$ be a polytope with $0 \in \operatorname{int}(K)$
given in $H$-representation $K = \{x : \langle a_i, x\rangle \leq 1\}_{i=1}^F$.
The \emph{polar body} (or \emph{dual}) is
\[
  K^\circ \coloneqq \{y \in \mathbb{R}^4 : \langle x, y \rangle \leq 1
    \text{ for all } x \in K\}
  = \operatorname{conv}(a_1, \ldots, a_F).
\]
\end{definition}
\end{unverified}

% Agent note: def:polar-body used in some .rs files is an alias
% for def:polytope-dual. The label below provides this as a cross-reference target.
\phantomsection\label{def:polar-body}% alias for cross-references from .rs files

\begin{unverified}
\begin{definition}[Cross product in $\mathbb{R}^4$]\label{def:cross-product-4d}
% [TODO: JÖRN - verify statement]
% Source: cross_product_4d.rs doc comment
The \emph{cross product} of three vectors $u, v, w \in \mathbb{R}^4$
is the unique vector $u \times v \times w \in \mathbb{R}^4$
perpendicular to all three, with magnitude equal to the volume
of the parallelepiped spanned by $u, v, w$, and orientation
determined by the right-hand rule.
Equivalently,
\[
  (u \times v \times w)_i
  = (-1)^{i+1} \det\bigl(M_i\bigr),
\]
where $M_i$ is the $3 \times 3$ matrix obtained by deleting
row~$i$ from the $4 \times 3$ matrix $(u \mid v \mid w)$.
\end{definition}
\end{unverified}

\subsection{Boundedness and exact vertex predicates}

\begin{unverified}
\begin{proposition}[Boundedness check]\label{prop:boundedness-iff}
Let $K = \{x : \langle a_\ell, x \rangle \le 1,\; \ell = 1, \ldots, F\}$
with $0 \in \operatorname{int} K$, and let $A$ be the $F \times 4$ matrix
with rows $a_1, \ldots, a_F$. Then $K$ is
bounded if and only if
\textup{(i)}~the dual vertices $a_1, \ldots, a_F$ span~$\mathbb{R}^4$, and
\textup{(ii)}~for every triple~$(i,j,k)$ whose dual vertices have a 1-dimensional
kernel direction~$d$, some $a_\ell$ has positive inner product with~$d$ and some
has negative inner product with~$d$.
% When rank(A) <= 2, no triple has a 1D kernel, so the triple check is
% vacuously true but K is unbounded. The "span R^4" clause handles this.
\end{proposition}

\begin{proof}
The recession cone $\operatorname{rec}(K) = \{d : \langle a_\ell, d \rangle \le 0\;\forall \ell\}$ satisfies
$\operatorname{rec}(K) \subset K$ (since $0 \in K$, any $d \in
\operatorname{rec}(K)$ gives a ray $\{td : t \ge 0\} \subset K$).
Conversely, if~$K$ is unbounded it contains a ray $\{x_0 + td : t \ge 0\}$
for some nonzero~$d$; substituting into $\langle a_\ell, x_0 + td \rangle \le 1$ and letting
$t \to \infty$ forces $\langle a_\ell, d \rangle \le 0$, so $d \in \operatorname{rec}(K)$.  Hence
\[
  K \text{ bounded}
  \;\Longleftrightarrow\;
  \operatorname{rec}(K) = \{0\}.
\]
Now $\operatorname{rec}(K) = \{0\}$ means there is no nonzero~$d$ with
$a_\ell \cdot d \le 0$ for all~$\ell$; equivalently, for every nonzero~$d$,
some $a_\ell \cdot d > 0$.  This is exactly the condition that the dual vertices
\emph{positively span}~$\mathbb{R}^4$.  It remains to show that positive span
is equivalent to the conjunction of~\textup{(i)} and~\textup{(ii)}.

\medskip\noindent\emph{($\Rightarrow$) Positive span implies both conditions.}
Positive span implies span~$\mathbb{R}^4$ (a fortiori). For the triple check:
let~$d$ be a kernel direction of triple~$(i,j,k)$, so $a_i \cdot d = a_j
\cdot d = a_k \cdot d = 0$. By positive span, some~$a_\ell$ has $a_\ell \cdot
d > 0$; since the triple's own dual vertices have zero dot product with~$d$, this
$a_\ell$ lies outside the triple. Applying the same argument to~$-d$ gives a
dual vertex with negative inner product.
% QC: The kernel is 1-dimensional by hypothesis, so d is determined up to
% sign. The condition "some positive, some negative" is sign-symmetric,
% hence independent of the choice of generator.

\medskip\noindent\emph{($\Leftarrow$) Span~$\mathbb{R}^4$ and all triples
pass implies positive span.}
We prove the contrapositive: if positive span fails but the dual vertices
span~$\mathbb{R}^4$, some triple fails.

Suppose $\operatorname{rec}(K) \neq \{0\}$, i.e., there exists a nonzero~$d$
with $a_\ell \cdot d \le 0$ for all~$\ell$. We produce a triple~$(i,j,k)$
whose kernel direction satisfies $a_\ell \cdot d \le 0$ for all~$\ell$
(so no dual vertex blocks it).

Among all unit vectors in~$\operatorname{rec}(K)$, pick~$d$ maximizing the
number of orthogonal dual vertices $S_0 = \{\ell : a_\ell \cdot d = 0\}$.  We claim
that~$\{a_\ell : \ell \in S_0\}$ contains three linearly independent vectors.
If not, $\{a_\ell : \ell \in S_0\}$ spans a subspace of dimension~$\le 2$,
so its orthogonal complement~$W$ has $\dim W \ge 4 - 2 = 2$.  Since
$a_\ell \cdot d = 0$ for every $\ell \in S_0$, we have $d \in W$.  Because
$\dim W \ge 2$ and $d \ne 0$, there exists $v \in W$ not parallel to~$d$.
Since $v \perp a_\ell$ for all $\ell \in S_0$, we have
\[
  a_\ell \cdot (d + tv) = a_\ell \cdot d + t\,(a_\ell \cdot v)
  = \begin{cases}
    0 + t \cdot 0 = 0 & \text{if } \ell \in S_0, \\
    \underbrace{a_\ell \cdot d}_{<\,0} + t\,(a_\ell \cdot v) < 0
    & \text{if } \ell \notin S_0, \text{ for small } |t|.
  \end{cases}
\]
% QC: ℓ ∉ S₀ is non-empty because dual vertices span R^4 ⟹ rank(A) = 4
% ⟹ ker(A) = {0}, so the nonzero vector d cannot satisfy a_ℓ · d = 0 for all ℓ.
(The case $\ell \notin S_0$ is non-empty: since the dual vertices span~$\mathbb{R}^4$,
the matrix~$A$ has trivial kernel, so~$d \ne 0$ cannot be orthogonal to all
dual vertices.)
So $d + tv \in \operatorname{rec}(K)$ for small~$|t|$.

Consider increasing~$|t|$ from zero.  The set
$T = \{t \in \mathbb{R} : a_\ell \cdot (d + tv) \le 0 \;\forall \ell\}$ is
a closed interval containing~$0$ (the intersection of half-lines
$\{t : a_\ell \cdot d + t\,a_\ell \cdot v \le 0\}$).  This interval has
at least one finite endpoint: if it extended indefinitely in both directions,
then $a_\ell\cdot v=0$ for every~$\ell$, contradicting
$\operatorname{rank}(A)=4$.  It contains a neighborhood of zero, so any
finite endpoint~$t^*$ is nonzero.  At $t = t^*$, some new index
$\ell^* \notin S_0$ satisfies $a_{\ell^*} \cdot (d + t^* v) = 0$.  Set
$d' = d + t^* v$; this is nonzero because $d$ and~$v$ are not parallel.
Since $v \in W$, we still have $a_\ell \cdot d' = a_\ell \cdot d = 0$
for every $\ell \in S_0$.  Thus
$\{\ell : a_\ell \cdot d' = 0\} \supseteq S_0 \cup \{\ell^*\}$,
a strict superset of~$S_0$.  Since $d' \in \operatorname{rec}(K)$, the unit
vector $d' / \|d'\|$ contradicts the choice of~$d$ as maximizing~$|S_0|$.

So~$\{a_\ell : \ell \in S_0\}$ contains three linearly independent vectors;
pick such a triple~$(i,j,k)$ from~$S_0$.  Since $a_i, a_j, a_k$ are linearly
independent in~$\mathbb{R}^4$, their common kernel $\ker(a_i, a_j, a_k)$ is
$(4 - 3)$-dimensional, so the triple falls under condition~\textup{(ii)}.
The direction~$d$ spans this kernel, and $d \in \operatorname{rec}(K)$ gives
$a_\ell \cdot d \le 0$ for all~$\ell$.  In particular, no dual vertex has positive
inner product with~$d$, so triple~$(i,j,k)$ fails condition~\textup{(ii)}.
\end{proof}
\end{unverified}

\begin{unverified}
\begin{lemma}[Positive span implies boundedness]\label{lem:positive-span}
% [TODO: JÖRN - verify statement and proof]
% Source: check_bounded_f64_first in vertex_enumeration/boundedness.rs
% Proof compressed here by reference to the preceding recession argument;
% this edited lemma remains unverified.
Let $a_1, \ldots, a_F \in \mathbb{R}^4 \setminus \{0\}$ and
$K = \{x \in \mathbb{R}^4 : \langle a_i, x \rangle \leq 1 \;\forall i\}$.
Then $K$ is bounded if and only if $\{a_1, \ldots, a_F\}$ positively
spans~$\mathbb{R}^4$, i.e., for every $d \in \mathbb{R}^4 \setminus \{0\}$
there exists $i$ with $\langle a_i, d \rangle > 0$.
\end{lemma}

\begin{proof}
The recession-cone argument in Proposition~\ref{prop:boundedness-iff} shows
that $K$ is bounded exactly when there is no nonzero $d$ with
$\langle a_i,d\rangle\leq 0$ for every~$i$.  This is precisely the stated
positive-spanning condition.
\end{proof}
\end{unverified}


\begin{unverified}
\begin{lemma}[Algorithmic bounded check via triples]\label{lem:bounded-triples}
% [TODO: JÖRN - verify statement and proof]
% Source: check_bounded_f64_first in vertex_enumeration/boundedness.rs
Under the hypotheses of Lemma~\ref{lem:positive-span}, assume
$\operatorname{rank}\{a_1, \ldots, a_F\} = 4$.
Then $\{a_i\}$ positively spans $\mathbb{R}^4$ if and only if for every
linearly independent triple $\{a_i,a_j,a_k\}$ and every
$d \in \ker(a_i, a_j, a_k) \setminus \{0\}$, there exist $l, m$ with
$\langle a_l, d \rangle > 0$ and $\langle a_m, d \rangle < 0$.
\end{lemma}

\begin{proof}
This is the triple part of Proposition~\ref{prop:boundedness-iff}, whose
recession-cone proof also covers redundant listed rows.  The rank assumption
and restriction to independent triples are essential: a dependent triple has
a kernel of dimension at least two, so there is no single triple-normal
direction to test.
\end{proof}
\end{unverified}


\begin{unverified}
\begin{lemma}[Exact vertex enumeration]\label{lem:vertex-enumeration}
% [TODO: JÖRN - verify statement and proof]
% Source: enumerate_vertices_int in vertex_enumeration/enumerate.rs
Let $K = \{x \in \mathbb{R}^4 : \langle a_i, x \rangle \leq 1 \;\forall i\}$
be a bounded polytope with $0 \in \operatorname{int}(K)$.
A point $v \in \mathbb{R}^4$ is a vertex of~$K$ if and only if there exists
a subset $S \subset \{1,\ldots,F\}$ with $|S| \geq 4$ such that:
\begin{enumerate}
  \item the matrix $(a_s)_{s \in S'}$ is nonsingular for some
    $S' \subset S$ with $|S'| = 4$,
  \item $\langle a_s, v \rangle = 1$ for all $s \in S$ (the defining
    constraints are tight), and
  \item $\langle a_i, v \rangle \leq 1$ for all $i \notin S$ (the
    remaining constraints are satisfied).
\end{enumerate}
The algorithm iterates over all $\binom{F}{4}$ four-element subsets $S'$,
solves $\langle a_{s_k}, v \rangle = 1$ ($k = 1,\ldots,4$) via Cramer's
rule, checks feasibility, and records the full incidence set
$S = \{i : \langle a_i, v \rangle = 1\}$.
\end{lemma}

\begin{proof}
($\Rightarrow$) Let $v$ be a vertex of~$K$. Then $v$ is an extreme point:
$v$ cannot be written as a convex combination of two distinct points in~$K$.
Let $S = \{i : \langle a_i, v \rangle = 1\}$ be the active set.
The active constraints must have $\operatorname{rank}\{a_s : s \in S\} = 4$,
since otherwise $v$ would not be an extreme point (one could perturb $v$
within the nullspace while remaining feasible). So $|S| \geq 4$ and some
4-element subset $S' \subset S$ has nonsingular coefficient matrix.
Feasibility of $v$ gives condition (3).

($\Leftarrow$) If conditions (1)--(3) hold, then $v$ is the unique solution
of the $4 \times 4$ nonsingular system $\langle a_{s_k}, v \rangle = 1$,
so $v$ cannot be written as a convex combination of two distinct feasible
points (they would also need to satisfy the 4 tight constraints, forcing
them to equal~$v$). Hence $v$ is an extreme point, i.e., a vertex.

The algorithm finds all vertices because every vertex has at least one
4-element subset of active constraints with nonsingular coefficient matrix
(by the rank argument above), and the algorithm tests all such subsets.
Deduplication handles non-simple vertices found via multiple subsets.
\end{proof}
\end{unverified}


\begin{unverified}
\begin{lemma}[Irredundancy check]\label{lem:irredundancy}
% [TODO: JÖRN - verify statement and proof]
% Source: check_irredundancy_f64_first in vertex_enumeration/irredundancy.rs
Let $K$ be a bounded, full-dimensional normalized $H$-polytope as in
Lemma~\ref{lem:vertex-enumeration}, with vertices $V$ and vertex-facet incidence
$V_i = \{v \in V : \langle a_i, v \rangle = 1\}$.
Assume $a_i\ne a_j$ for every $j\ne i$.  Then the $i$-th listed constraint
is irredundant if and only if
$\operatorname{aff}(V_i)$ is a hyperplane, equivalently
$\operatorname{affrank}(V_i) = 3$.
Without the uniqueness assumption, affine rank three still says that the
supporting hyperplane defines a facet, but duplicate copies of its inequality
are individually redundant.
\end{lemma}

\begin{proof}
The $i$-th constraint defines the face
$F_i = K \cap \{x : \langle a_i, x \rangle = 1\}$.
The vertices of $F_i$ are exactly $V_i$ (the vertices of~$K$ that lie
on the hyperplane $\langle a_i, x \rangle = 1$).

If $\operatorname{affrank}(V_i)=3$, choose a point in the relative interior
of $F_i$ away from the finitely many intersections with distinct listed
hyperplanes.  Every other listed inequality is strict there, so a small
outward displacement violates only the $i$-th inequality.  Removing it
strictly enlarges~$K$.

Conversely, if removing the $i$-th inequality enlarges~$K$, take a point
that satisfies all other inequalities and violates the $i$-th.  A small
convex combination with the origin still violates the $i$-th and makes
every other inequality strict.  The segment from the origin to this point
crosses the $i$-th supporting hyperplane where all other inequalities are
strict, giving a relatively open three-dimensional patch of $F_i$.  Thus
$\operatorname{affrank}(V_i)=3$.
\end{proof}
\end{unverified}

The affine-rank check alone certifies a facet-supporting row.  A validator
that admits duplicate listed rows cannot infer irredundancy of each listed
inequality from that check alone.


\begin{unverified}
\begin{proposition}[Integer-scaled Cramer's rule]\label{prop:integer-cramer}
% [TODO: JÖRN - verify statement and proof]
% Source: optimized vertex enumeration in vertex_enumeration/enumerate.rs
Let $a_1, \ldots, a_F \in \mathbb{Q}^4$ be the dual vertices defining
$K = \{x \in \mathbb{R}^4 : \langle a_i, x \rangle \leq 1\}$.
Let $D = \operatorname{lcm}$ of all denominators of all components of
$a_1, \ldots, a_F$, and set $A_i = D \cdot a_i \in \mathbb{Z}^4$.

For a four-element subset $S = \{s_0, s_1, s_2, s_3\}$, let $M_S$ be the
$4 \times 4$ matrix with rows $A_{s_k}$, and let $\delta = \det(M_S)$.
If $\delta \neq 0$, the candidate vertex~$v$ solving
$\langle a_{s_k}, v \rangle = 1$ for all $k$ has coordinates
\[
  v^{(j)} = \frac{D \cdot \nu_j}{\delta},
  \qquad
  \nu_j = \det\!\bigl(M_S \text{ with column } j
    \text{ replaced by } \mathbf{1}\bigr),
\]
where $\mathbf{1} = (1,1,1,1)^\top$.

The feasibility gap for facet $i \notin S$ is
\[
  1 - \langle a_i, v \rangle
  = \frac{g_i}{\delta},
  \qquad
  g_i = \delta - \sum_{j=0}^{3} A_i^{(j)} \nu_j \;\in\; \mathbb{Z}.
\]
Hence $1 - \langle a_i, v \rangle \geq 0$ if and only if
$\operatorname{sign}(g_i) = \operatorname{sign}(\delta)$ or $g_i = 0$,
and all intermediate computations are in~$\mathbb{Z}$.
\end{proposition}

\begin{proof}
The system $\langle a_{s_k}, v \rangle = 1$ for $k = 0, \ldots, 3$
is equivalent to $M_S\,v = D\mathbf{1}$.
By Cramer's rule,
\[
  v^{(j)} = \frac{\det(M_S^{(j)})}{\delta},
\]
where $M_S^{(j)}$ is $M_S$ with column~$j$ replaced by
$D \cdot \mathbf{1}$.
By multilinearity of the determinant in the $j$-th column,
$\det(M_S^{(j)}) = D \cdot \nu_j$.

For the feasibility gap:
\[
  1 - \langle a_i, v \rangle
  = 1 - \sum_{j=0}^{3} \frac{A_i^{(j)}}{D} \cdot \frac{D\, \nu_j}{\delta}
  = 1 - \frac{1}{\delta} \sum_{j=0}^{3} A_i^{(j)} \nu_j
  = \frac{\delta - \sum_j A_i^{(j)} \nu_j}{\delta}
  = \frac{g_i}{\delta}.
\]
Since $A_i^{(j)}, \nu_j, \delta \in \mathbb{Z}$, we have $g_i \in \mathbb{Z}$.
The sign condition follows: $g_i / \delta \geq 0$ iff
$\operatorname{sign}(g_i) = \operatorname{sign}(\delta)$ or $g_i = 0$.
\end{proof}
\end{unverified}

% Moved from formal/rational-integer-scaling.tex; exact integer scaling for
% rational polytope predicates.
%
% Labels defined here:
%   lem:rational-common-denominator-signs,
%   lem:integer-scaled-origin-interior-test,
%   lem:integer-scaled-polar-candidates,
%   lem:integer-scaled-feasibility-gap
%
% Status: mechanical agent-written proof note, unapproved by Jorn. It records
% the proof obligation used by the BigRational specializations in
% crates/euclidean-polytopes/src/predicates.rs and
% crates/euclidean-polytopes/src/polar.rs.

\subsection{Common-denominator scaling}

Let $v_1,\ldots,v_N \in \mathbb{Q}^4$.  Choose any positive integer $D$
such that $p_i := Dv_i \in \mathbb{Z}^4$ for all $i$.  These lemmas need only
the common-denominator property; the referenced Rust paths choose the least
common multiple of the component denominators.

\begin{lemma}[Common-denominator scaling preserves signs]
\label{lem:rational-common-denominator-signs}
For every $x \in \mathbb{Q}^4$ and every integer $D > 0$,
\[
  \operatorname{sign}\langle v_i,x\rangle
  =
  \operatorname{sign}\langle Dv_i,x\rangle .
\]
Moreover, for any three indices $i,j,k$, the $4$-dimensional cross product
normal satisfies
\[
  n(Dv_i,Dv_j,Dv_k) = D^3 n(v_i,v_j,v_k),
\]
so its dot products against scaled points satisfy
\[
  \langle Dv_\ell,n(Dv_i,Dv_j,Dv_k)\rangle
  =
  D^4\langle v_\ell,n(v_i,v_j,v_k)\rangle .
\]
\end{lemma}

\begin{proof}
The first identity is multiplication by the positive scalar $D$.  The cross
product formula is trilinear in its three arguments, hence contributes a factor
$D^3$.  The final identity adds the fourth factor $D$ from the tested point.
All factors are positive, so signs are unchanged.
\end{proof}

\begin{lemma}[Integer-scaled origin-interior test]
\label{lem:integer-scaled-origin-interior-test}
The rational triple-normal positive-spanning test for
$0 \in \operatorname{int}\operatorname{conv}\{v_i\}$ returns the same answer
after replacing all $v_i$ by $p_i = Dv_i$.
\end{lemma}

\begin{proof}
The rank precondition is unchanged by multiplication by the invertible scalar
$D$.  For every tested triple, Lemma~\ref{lem:rational-common-denominator-signs}
shows that the normal is zero before scaling if and only if it is zero after
scaling, and that every point-normal dot-product sign is unchanged.  Therefore
each separating-triple clause has the same truth value in the rational and
integer-scaled tests.
\end{proof}

\begin{lemma}[Integer-scaled polar candidates]
\label{lem:integer-scaled-polar-candidates}
Let $A$ be the $4 \times 4$ rational matrix with rows $v_{i_1},\ldots,v_{i_4}$,
and let $P = DA$.  If $\det P \neq 0$, then the exact solution of
$Ay=\mathbf{1}$ is recovered by solving $Pz=D\mathbf{1}$ and setting $y=z$.
Equivalently, Cramer's rule may be evaluated over integers using $P$ and the
right-hand side $D\mathbf{1}$.
\end{lemma}

\begin{proof}
Since $P=DA$, the equation $Ay=\mathbf{1}$ is equivalent to $Py=D\mathbf{1}$.  Cramer's
rule for $Py=D\mathbf{1}$ uses only integer determinants.  The implementation
chooses the least common multiple for $D$, while the argument holds for any
positive common denominator.  Thus determinant zeros are exact rational
ones, not rounded approximations.
\end{proof}

\begin{lemma}[Integer-scaled feasibility and incidence gap]
\label{lem:integer-scaled-feasibility-gap}
Let $v=p/D$ with $p\in\mathbb{Z}^4$ and $D>0$.  Let a rational candidate be
stored as $y=n/d$, where $n\in\mathbb{Z}^4$ and $d>0$.  Then
\[
  G(p,n,d,D) := Dd-\langle p,n\rangle
\]
has the same sign as $1-\langle v,y\rangle$.  In particular,
$G(p,n,d,D) > 0$ is exact strict feasibility for that row,
$G(p,n,d,D) = 0$ is exact incidence, and $G(p,n,d,D) < 0$ is exact violation
of the half-space inequality $\langle v,y\rangle \le 1$.
\end{lemma}

\begin{proof}
By substituting $v=p/D$ and $y=n/d$,
\[
  1-\langle v,y\rangle
  =
  1-\frac{\langle p,n\rangle}{Dd}
  =
  \frac{Dd-\langle p,n\rangle}{Dd}.
\]
The denominator $Dd$ is positive, so the numerator and the rational gap have
the same sign.  The three sign cases are exactly the half-space interior,
boundary, and exterior cases.
\end{proof}
